%%
%% CHI Zen starter paper.
%%
%% Compile with LuaLaTeX (Overleaf: Menu > Settings > Compiler > LuaLaTeX,
%% TeX Live version 2026 or later). See README.md for details.
%%
%% \DocumentMetadata must come before \documentclass. It switches on a
%% tagged (accessible) PDF that aims at PDF/UA-2, plus PDF/A-4f for
%% archiving, as recommended by the acmart documentation (v2.17+).
%% If you are stuck on an old TeX Live, comment it out: the paper still
%% compiles, but the PDF will not be tagged.
\DocumentMetadata{
  lang          = en,
  pdfstandard   = ua-2,
  pdfstandard   = a-4f,
  tagging       = on,
  tagging-setup = {math/setup=mathml-SE},
}

%% Class options for CHI (from the CHI 2027 "CHI Publication Formats" page,
%% https://chi2027.acm.org/chi-publication-formats/):
%%   anonymous submission:      \documentclass[manuscript,review,anonymous]{acmart}
%%   non-anonymous submission:  \documentclass[manuscript,review]{acmart}
%%   publication (camera-ready, after acceptance): \documentclass[sigconf]{acmart}
%% Always check the current year's call for papers before you submit.
%% natbib=false is ours, not CHI's: this template uses biblatex (see below).
\documentclass[manuscript,review,anonymous,natbib=false]{acmart}

% Turn on the CHIZen features
%
%  This file imports each of the separate .tex files in this directory -- feel free to enable/disable each separately as you'd like
%



%
%   A bunch of common includes for text formatting, features. 
%
% Accented characters: \usepackage[utf8]{inputenc} is no longer needed.
% UTF-8 has been the default input encoding of LaTeX since 2018, and
% LuaLaTeX reads UTF-8 natively.

% CJK (Chinese, Japanese, Korean) characters: \cjk{改善}
% CJKutf8 only works with pdfLaTeX, so under LuaLaTeX (the engine we use
% for tagged, accessible PDFs) CJK text is set in the Harano Aji Mincho
% font that ships with TeX Live (and therefore with Overleaf).
% Final mode (TAPS) loads neither iftex nor fontspec (acmart loads both
% under LuaLaTeX) nor CJKutf8 (not on ACM's list), so there \cjk needs
% LuaLaTeX.
\ifchizenfinal
\ifLuaTeX
  \newfontfamily\chizenCJKfont{HaranoAjiMincho-Regular.otf}
  \newcommand{\cjk}[1]{{\chizenCJKfont #1}}
\else
  \newcommand{\cjk}[1]{\PackageError{CHIZen}{\string\cjk\space needs LuaLaTeX in final mode}{%
    CJKutf8 is not on ACM's TAPS list. Compile with LuaLaTeX.}}
\fi
\else
\usepackage{iftex}
\ifLuaTeX
  \usepackage{fontspec}
  \newfontfamily\chizenCJKfont{HaranoAjiMincho-Regular.otf}
  \newcommand{\cjk}[1]{{\chizenCJKfont #1}}
\else
  \usepackage{CJKutf8}
  \newcommand{\cjk}[1]{\begin{CJK}{UTF8}{min}#1\end{CJK}}
\fi
\fi


% ACCESSIBILITY
% The old `accessibility' package is obsolete and must not be loaded.

% MARKDOWN (draft mode only: ACM's TAPS does not accept it)
%
% TODO: Explore how useful this is for general-purpose HCI papers
%       - Tables/Figures need to be a separate thing, but useful for everything else?
%
% markdown v3 loads amssymb, which clashes with the newtxmath fonts that
% acmart already loads (both define \Bbbk etc.). newtxmath provides all the
% amssymb symbols, so tell LaTeX amssymb is already there.
\ifchizenfinal\else
\makeatletter
\@ifpackageloaded{newtxmath}{\@namedef{ver@amssymb.sty}{2013/01/14 v3.01 (symbols provided by newtxmath)}}{}
\makeatother
% `notes' was called `footnotes' before markdown v3.
% LaTeX inside a markdown block: maths between dollar signs (texMathDollars),
% any other command as a raw span, `\citep{key}`{=tex} (rawAttribute), and
% % comments are dropped (texComments). These replace the `hybrid' option,
% which markdown 3.7.1 soft-deprecated. Plain text is now escaped, so write
% & as it is (not \\&); a % still starts a comment, so write \% for a percent.
\usepackage[notes,definitionLists,hashEnumerators,smartEllipses,rawAttribute,texMathDollars,texComments]{markdown}

\usepackage{paracol}
\usepackage{minted}
% fvextra (under minted) moves each code line up by 0.25pt with \vspace, so
% that background colours of adjacent lines overlap; acmart warns about any
% \vspace in the text. Listings here have no background colour, so turn it
% off. (With bgcolor=..., add bgcolorboxoverlap=0.25pt to that listing.)
\fvset{backgroundcolorboxoverlap=0pt}
% acmart's monospaced font, Inconsolata, has no italic, which minted styles
% use for comments. Say so (set them upright, as LaTeX would after a "Font
% shape undefined" warning). No spaces in the size specification: it is
% read from a macro.
\makeatletter
\begingroup\ttfamily\endgroup % declares the font family
\ifcsname\encodingdefault/\ttdefault/m/it\endcsname\else
  \edef\chizen@tmp{\noexpand\DeclareFontShape{\encodingdefault}{\ttdefault}{m}{it}{<->ssub*\ttdefault/m/n}{}}%
  \chizen@tmp
\fi
\makeatother
\fi
\ifchizenfinal
% Final mode: markdown, paracol and minted are not on ACM's TAPS list.
% Using them stops with an error; the body of the environment is skipped
% (with the comment package, which acmart loads, so verbatim text is fine).
\makeatletter
\def\chizen@textnotinfinal#1#2{%
  \PackageError{CHIZen}{#1 is not available in final mode}{%
    ACM's TAPS does not accept #1, so CHI Zen does not load it when the
    paper says \string\def\string\chizenmode{final}. #2 See the CHI Zen
    manual, Section "Camera-Ready for ACM TAPS".}}
\excludecomment{markdown}
\AddToHook{env/markdown/begin}{\chizen@textnotinfinal{markdown}{Rewrite
  the Markdown block in LaTeX.}}
\NewDocumentCommand\markdownInput{O{}m}{\chizen@textnotinfinal{markdown}{Rewrite
  #2 in LaTeX.}}
\excludecomment{minted}
\AddToHook{env/minted/begin}{\chizen@textnotinfinal{minted}{Use the
  listings package (\string\usepackage{listings} and a lstlisting
  environment) instead.}}
\NewDocumentCommand\mint{}{\chizen@textnotinfinal{minted (\string\mint)}{Use
  \string\lstinline\space from the listings package instead.}}
\NewDocumentCommand\mintinline{}{\chizen@textnotinfinal{minted (\string\mintinline)}{Use
  \string\lstinline\space from the listings package instead.}}
\NewDocumentCommand\inputminted{O{}mm}{\chizen@textnotinfinal{minted (\string\inputminted)}{Use
  \string\lstinputlisting\space from the listings package instead.}}
\NewDocumentCommand\usemintedstyle{O{}m}{}
\NewDocumentEnvironment{paracol}{m +b}{\chizen@textnotinfinal{paracol}{Use
  a table or acmart's two-column layout instead.}}{}
\makeatother
\fi
\usepackage{fontawesome5}

% This package is useful for revisions! 
% It defines \comment, which acmart's comment package defines too: keep
% comment's \comment and rename changes' one to \SCcomment.
\ifchizenfinal
% Final mode: savesym is not on ACM's TAPS list, so do the same by hand,
% and accept every change ([final]: no colours or struck-out text in the
% camera-ready version). commentmarkup=footnote keeps changes from loading
% todonotes, which loads TikZ.
\let\chizenSavedComment\comment
\let\comment\undefined
\usepackage[final,commentmarkup=footnote]{changes}
\let\SCcomment\comment
\let\comment\chizenSavedComment
\else
\usepackage{savesym}
\savesymbol{comment}
\let\comment\undefined % savesym leaves \comment as \relax, which changes v4 treats as defined
\usepackage{changes}
%\usepackage[final]{changes}
\restoresymbol{SC}{comment}
\fi


% Mark likes to add these options to his papers .... worth considering
\frenchspacing
\raggedbottom

% Lennart's acronym package
\usepackage[nolist]{acronym}

\begin{acronym}
\acro{ANOVA}{Analysis of Variance}
\acro{HCI}{Human-Computer Interaction}
\end{acronym}


% Jim's quotation stuff
%\renewenvironment{quote}{%
%  \list{}{%
%    \leftmargin0.4cm   % this is the adjusting screw
%    \rightmargin\leftmargin
%  }
%  \item\relax
%}
%{\endlist}

% Jim added this -- useful highlighting function \hl{}
%\usepackage{soul,color}



% Study Variables
% ----------------------------

% scale up the sc font slightly
\ifchizenfinal
% Final mode: letltxmacro, lipsum and relsize are not on ACM's TAPS list.
% \NewCommandCopy is built into LaTeX; \chizenrelscale does what
% \relscale does.
\makeatletter
\NewCommandCopy{\oldscshape}{\scshape}
\def\chizenrelscale#1{%
  \@tempdima=#1\dimexpr\f@size pt\relax
  \@tempdimb=#1\dimexpr\f@baselineskip\relax
  \edef\chizen@tmp{\noexpand\fontsize{\the\@tempdima}{\the\@tempdimb}}%
  \chizen@tmp\selectfont}
\renewcommand{\scshape}[1]{\oldscshape\chizenrelscale{1.2}#1}
\NewDocumentCommand\lipsum{O{}}{\PackageError{CHIZen}{\string\lipsum\space
  is not available in final mode}{Replace the placeholder text with your own.}}
\makeatother
\else
\usepackage{letltxmacro,lipsum,relsize}
\LetLtxMacro{\oldscshape}{\scshape}
\renewcommand{\scshape}[1]{\oldscshape\relscale{1.2}#1}
\fi

% an experiment factor/variable
\newcommand{\f}[1]{\mbox{{\textsc{#1}}}\xspace}
% usage: \f{block}, \f{width}, etc.
% NOTE: need \xspace package to fix spacing issues
%       \mbox prevents breaking

% a code when qualitative coding 
\newcommand{\cd}[1]{\textit{``{#1}''}\xspace}
% usage: \cd{errors}, \cd{work}


%
%   Figures -- formerly IDEA Plots -- Wraps most features from PGFPlots
%       * Handles common formatting, defines some common graph types
%       * See file for detailed documentation
%
%%%%%%%%%%%%%%%%%%%%%%%%%%%%%%%%%%%%%%%%%%%%%%%%%
%  HOW TO USE
%%%%%%%%%%%%%%%%%%%%%%%%%%%%%%%%%%%%%%%%%%%%%%%%%
%
% CHIZen/init.tex loads this file. The rules for choosing and using each
% style (and the matplotlib/ggplot equivalents) are in the CHI Zen guide:
%   plugins/chizen/skills/chizen/references/figures.md
%
% Add a graph to your paper with this template:
%
%\begin{tikzpicture}
%\begin{axis}[<STYLE>,
%	% X Axis options (xlabel, xmin, xmax, xtick={})
%	% Y Axis options (ylabel, ymin, ymax, ytick={})
%       ]
%
%    % DATA GOES IN HERE (preferably read from data/*.csv)
%
%\end{axis}
%\end{tikzpicture}
%
%
% OR, for multiple plots use this template:
%
%\begin{tikzpicture}
%\begin{groupplot}[
%   group style={
%       group size= <COL> by <ROW>,
%       x descriptions at=edge bottom,
%       y descriptions at=edge left,
%       vertical sep=0pt,
%       horizontal sep=0pt},
%  <STYLE>,
%	% X Axis options (xmin, xmax, xtick={})
%	% Y Axis options (ymin, ymax, ytick={})
%      ]
%
%    % DATA GOES IN HERE
%
%\end{groupplot}
%\end{tikzpicture}
%
%
% The following styles are defined in this library:
%
%   chizen bar          clustered bar graph (comparing conditions)
%   chizen stacked bar  stacked bar graph
%   chizen box          box plot, filled boxes with a white median (distributions)
%   chizen tufte box    Tufte's minimal box plot (a dot for the median, no box)
%   chizen tufte panel  bare-bones bar panel for Tufte-style stacked panels (groupplot)
%   chizen line         line graph (trends over time)
%   chizen Likert       diverging stacked bar for Likert responses
%   range frame         Tufte's range frame (axis lines span the data only)
% and the cycle lists (palettes) chizen lines, okabe-ito, waterloo-light,
% waterloo-dark, colorbrewer-RYB and colorbrewer-RYB-plain.
%
% NAMES. The styles used to be called `IDEA ...' (after the IDEA lab), e.g.
% `IDEA bar' and the `IDEA lines' cycle list. Those names still work as
% aliases (see the end of this file), so existing papers compile unchanged,
% but new figures should use the `chizen ...' names.
%




%%%%%%%%%%%%%%%%%%%%%%%%%%%%%%%%%%%%%%%%%%%%%%%%%
%   Common Includes
%%%%%%%%%%%%%%%%%%%%%%%%%%%%%%%%%%%%%%%%%%%%%%%%%
\usepackage{pgfplots}
\usepackage{pgfplotstable}
\usepgfplotslibrary{dateplot}
\usepgfplotslibrary{groupplots}
\usepackage{pgfgantt}

\usepackage{subcaption}

\usetikzlibrary{external}
%\tikzexternalize[prefix=tikz/] % Must be in main .tex file?

\usetikzlibrary{pgfplots.statistics}
\usetikzlibrary{matrix}

\usepackage{bbding}

\usepackage{svg}

\pgfplotsset{compat=newest}

%%%%%%%%%%%%%%%%%%%%%%%%%%%%%%%%%%%%%%%%%%%%%%%%%
%   Colour Definitions
%%%%%%%%%%%%%%%%%%%%%%%%%%%%%%%%%%%%%%%%%%%%%%%%%

\definecolor{RYB1}{RGB}{141, 211, 199}
\definecolor{RYB2}{RGB}{255, 255, 179} % NOTE: near-invisible on white (1.04:1); use colorbrewer-RYB (with outlines), not -plain
\definecolor{RYB3}{RGB}{190, 186, 218}
\definecolor{RYB4}{RGB}{251, 128, 114}
\definecolor{RYB5}{RGB}{128, 177, 211}
\definecolor{RYB6}{RGB}{253, 180, 98}
\definecolor{RYB7}{RGB}{179, 222, 105}

\pgfplotscreateplotcyclelist{colorbrewer-RYB}{
{RYB1!50!black,fill=RYB1},
{RYB2!50!black,fill=RYB2},
{RYB3!50!black,fill=RYB3},
{RYB4!50!black,fill=RYB4},
{RYB5!50!black,fill=RYB5},
{RYB6!50!black,fill=RYB6},
{RYB7!50!black,fill=RYB7},
}

\pgfplotscreateplotcyclelist{colorbrewer-RYB-plain}{
{RYB1},
{RYB2},
{RYB3},
{RYB4},
{RYB5},
{RYB6},
{RYB7},
}

%%%%%%%%%%%%%%%%%%%%
% WATERLOO PALETTES - DIGITAL
%%%%%%%%%%%%%%%%%%%%

\definecolor{AHSLight}{RGB}{0, 154, 166}
\definecolor{AHSDark}{RGB}{0, 127, 138}

\definecolor{ArtsLight}{RGB}{233, 131, 0}
\definecolor{ArtsDark}{RGB}{172, 97, 0}

\definecolor{EngineeringLight}{RGB}{204, 170, 255}
\definecolor{EngineeringDark}{RGB}{87, 6, 140}

\definecolor{EnvironmentLight}{RGB}{182, 191, 0}
\definecolor{EnvironmentDark}{RGB}{116, 120, 0}

\definecolor{MathLight}{RGB}{255, 136, 221}
\definecolor{MathDark}{RGB}{224, 36, 154}
 
\definecolor{ScienceLight}{RGB}{119, 187, 225}
\definecolor{ScienceDark}{RGB}{0, 115, 207}


\definecolor{SchoolRedLight}{RGB}{247, 119, 119}
\definecolor{SchoolRedDark}{RGB}{150, 23, 46}



\pgfplotscreateplotcyclelist{waterloo-light}{
{AHSLight!50!black,fill=AHSLight},
{ArtsLight!50!black,fill=ArtsLight},
{EngineeringLight!50!black,fill=EngineeringLight},
{EnvironmentLight!50!black,fill=EnvironmentLight},
{MathLight!50!black,fill=MathLight},
{ScienceLight!50!black,fill=ScienceLight},
{SchoolRedLight!50!black,fill=SchoolRedLight},
}

\pgfplotscreateplotcyclelist{waterloo-dark}{
{AHSDark!50!black,fill=AHSDark},
{ArtsDark!50!black,fill=ArtsDark},
{EngineeringDark!50!black,fill=EngineeringDark},
{EnvironmentDark!50!black,fill=EnvironmentDark},
{MathDark!50!black,fill=MathDark},
{ScienceDark!50!black,fill=ScienceDark},
{SchoolRedDark!50!black,fill=SchoolRedDark},
}


%%%%%%%%%%%%%%%%%%%%
% GREYSCALE / COLOUR-BLIND NOTES (relative luminance, WCAG 2 formula)
%   None of the categorical palettes above is greyscale-safe: some colours
%   have almost the same luminance and print as the same grey, e.g.
%     waterloo-light: ArtsLight .336 vs SchoolRedLight .343,
%                     MathLight .441 vs ScienceLight .449
%     waterloo-dark:  AHS, Arts, Environment, Science and Math are all .17-.19
%     colorbrewer-RYB: RYB1 .564 vs RYB6 .544; RYB2 (.960) is almost white.
%   The colours are kept as they are because published figures use them.
%   For colour-blind readers use okabe-ito below; for black-and-white print,
%   also vary marks, dash patterns or hatching, or label series directly.
%%%%%%%%%%%%%%%%%%%%

%%%%%%%%%%%%%%%%%%%%
% OKABE-ITO PALETTE - colour-blind safe (Okabe & Ito, "Color Universal Design")
%   Distinguishable with the common forms of colour blindness, but NOT in
%   greyscale (OkabeItoOrange .416 vs OkabeItoSkyBlue .405). OkabeItoYellow
%   is light on white (1.32:1), so the cycle list gives every fill a darker
%   outline, like the palettes above.
%%%%%%%%%%%%%%%%%%%%

\definecolor{OkabeItoOrange}{RGB}{230, 159, 0}
\definecolor{OkabeItoSkyBlue}{RGB}{86, 180, 233}
\definecolor{OkabeItoBluishGreen}{RGB}{0, 158, 115}
\definecolor{OkabeItoYellow}{RGB}{240, 228, 66}
\definecolor{OkabeItoBlue}{RGB}{0, 114, 178}
\definecolor{OkabeItoVermillion}{RGB}{213, 94, 0}
\definecolor{OkabeItoReddishPurple}{RGB}{204, 121, 167}
\definecolor{OkabeItoBlack}{RGB}{0, 0, 0}

\pgfplotscreateplotcyclelist{okabe-ito}{
{OkabeItoOrange!50!black,fill=OkabeItoOrange},
{OkabeItoSkyBlue!50!black,fill=OkabeItoSkyBlue},
{OkabeItoBluishGreen!50!black,fill=OkabeItoBluishGreen},
{OkabeItoYellow!50!black,fill=OkabeItoYellow},
{OkabeItoBlue!50!black,fill=OkabeItoBlue},
{OkabeItoVermillion!50!black,fill=OkabeItoVermillion},
{OkabeItoReddishPurple!50!black,fill=OkabeItoReddishPurple},
{OkabeItoBlack,fill=OkabeItoBlack},
}






%%%%%%%%%%%%%%%%%%%%
% 5-Point DIVERGENT PALETTES - From ColorBrewer2
%%%%%%%%%%%%%%%%%%%%

\definecolor{Likert5_SD}{RGB}{166,97,26}
\definecolor{Likert5_D}{RGB}{223,194,125}
\definecolor{Likert5_N}{RGB}{245,245,245}
\definecolor{Likert5_A}{RGB}{128,205,193}
\definecolor{Likert5_SA}{RGB}{1,133,113}

\definecolor{Likert5_1_SD}{RGB}{166,97,26}
\definecolor{Likert5_1_D}{RGB}{223,194,125}
\definecolor{Likert5_1_N}{RGB}{245,245,245}
\definecolor{Likert5_1_A}{RGB}{128,205,193}
\definecolor{Likert5_1_SA}{RGB}{1,133,113}

\definecolor{Likert5_2_SD}{RGB}{123,50,148}
\definecolor{Likert5_2_D}{RGB}{194,165,207}
\definecolor{Likert5_2_N}{RGB}{247,247,247}
\definecolor{Likert5_2_A}{RGB}{166,219,160}
\definecolor{Likert5_2_SA}{RGB}{0,136,55}



%%%%%%%%%%%%%%%%%%%%
% 7-Point DIVERGENT PALETTES - From ColorBrewer2
%%%%%%%%%%%%%%%%%%%%

% Purple -> Green
\definecolor{Likert7_0}{RGB}{118,42,131}
\definecolor{Likert7_1}{RGB}{175,141,195}
\definecolor{Likert7_2}{RGB}{231,212,232}
\definecolor{Likert7_3}{RGB}{229,229,229}
\definecolor{Likert7_4}{RGB}{217,240,211}
\definecolor{Likert7_5}{RGB}{127,191,123}
\definecolor{Likert7_6}{RGB}{27,120,55}

% Orange/Brown -> Blue
\definecolor{Likert7_2_0}{RGB}{140,81,10}
\definecolor{Likert7_2_1}{RGB}{216,179,101}
\definecolor{Likert7_2_2}{RGB}{246,232,195}
\definecolor{Likert7_2_3}{RGB}{229,229,229}
\definecolor{Likert7_2_4}{RGB}{199,234,229}
\definecolor{Likert7_2_5}{RGB}{90,180,172}
\definecolor{Likert7_2_6}{RGB}{1,102,94}

% Orange -> Purple
\definecolor{Likert7_3_0}{RGB}{179,88,6}
\definecolor{Likert7_3_1}{RGB}{241,163,64}
\definecolor{Likert7_3_2}{RGB}{254,224,182}
\definecolor{Likert7_3_3}{RGB}{229,229,229}
\definecolor{Likert7_3_4}{RGB}{216,218,235}
\definecolor{Likert7_3_5}{RGB}{153,142,195}
\definecolor{Likert7_3_6}{RGB}{84,39,136}

% Blue -> Red
\definecolor{Likert7_4_0}{RGB}{33,102,172}
\definecolor{Likert7_4_1}{RGB}{103,169,207}
\definecolor{Likert7_4_2}{RGB}{209,229,240}
\definecolor{Likert7_4_3}{RGB}{229,229,229}
\definecolor{Likert7_4_4}{RGB}{253,219,199}
\definecolor{Likert7_4_5}{RGB}{239,138,98}
\definecolor{Likert7_4_6}{RGB}{178,24,43}


%%%%%%%%%%%%%%%%%%%%
% 7-Point DIVERGENT PALETTES - Jim for Avatar Paper
%%%%%%%%%%%%%%%%%%%%

% Blue -> Green to match Cleric
\definecolor{Cleric_0}{RGB}{33,102,172}
\definecolor{Cleric_1}{RGB}{103,169,207}
\definecolor{Cleric_2}{RGB}{209,229,240}
\definecolor{Cleric_3}{RGB}{229,229,229}
\definecolor{Cleric_4}{RGB}{217,240,211}
\definecolor{Cleric_5}{RGB}{127,191,123}
\definecolor{Cleric_6}{RGB}{27,120,55}

% Orange -> Red
\definecolor{Monster_0}{RGB}{178,24,43}
\definecolor{Monster_1}{RGB}{239,138,98}
\definecolor{Monster_2}{RGB}{253,219,199}
\definecolor{Monster_3}{RGB}{229,229,229}
\definecolor{Monster_4}{RGB}{216,218,235}
\definecolor{Monster_5}{RGB}{153,142,195}
\definecolor{Monster_6}{RGB}{84,39,136}


\pgfplotscreateplotcyclelist{avatar}{Cleric_0, Monster_0}

%%%%%%%%%%%%%%%%%%%%
% LINE GRAPHS - default for chizen line
%   Each series differs in colour, dash pattern AND mark, so lines can be told
%   apart when printed in black and white. Colours are the Okabe-Ito palette
%   (colour-blind safe), ordered darkest first so thin lines stay visible;
%   OkabeItoYellow is left out because a thin yellow line on white is too
%   faint (1.32:1). Marks are always drawn solid. For dense data (e.g.,
%   monthly counts) add `no markers' to the axis and rely on the dash patterns.
%%%%%%%%%%%%%%%%%%%%

\pgfplotscreateplotcyclelist{chizen lines}{
{OkabeItoBlue, solid, mark=*, mark options={solid}},
{OkabeItoVermillion, dashed, mark=square*, mark options={solid}},
{OkabeItoBluishGreen, dotted, mark=triangle*, mark options={solid}},
{OkabeItoReddishPurple, dashdotted, mark=diamond*, mark options={solid}},
{OkabeItoOrange, densely dashed, mark=o, mark options={solid}},
{OkabeItoSkyBlue, densely dotted, mark=square, mark options={solid}},
{OkabeItoBlack, loosely dashed, mark=triangle, mark options={solid}},
}

%%%%%%%%%%%%%%%%%%%%%%%%%%%%%%%%%%%%%%%%%%%%%%%%%
%   Outer South Legend Hack
%%%%%%%%%%%%%%%%%%%%%%%%%%%%%%%%%%%%%%%%%%%%%%%%%
\makeatletter
\pgfplotsset{
    every axis x label/.append style={
        alias=current axis xlabel
    },
    legend pos/outer south/.style={
        /pgfplots/legend style={
            at={%
                (%
                \@ifundefined{pgf@sh@ns@current axis xlabel}%
                {xticklabel cs:0.5}%
                {current axis xlabel.south}%
                )%
            },
            anchor=north
        }
    }
}

%%%%%%%%%%%%%%%%%%%%%%%%%%%%%%%%%%%%%%%%%%%%%%%%%
%   Grid Figures Label Hack
%%%%%%%%%%%%%%%%%%%%%%%%%%%%%%%%%%%%%%%%%%%%%%%%%

\pgfplotsset{
    groupplot xlabel/.initial={},
    every groupplot x label/.style={
        at={($({\pgfplots@group@name\space c1r\pgfplots@group@rows.west}|-{\pgfplots@group@name\space c1r\pgfplots@group@rows.outer south})!0.5!({\pgfplots@group@name\space c\pgfplots@group@columns r\pgfplots@group@rows.east}|-{\pgfplots@group@name\space c\pgfplots@group@columns r\pgfplots@group@rows.outer south})$)},
        anchor=north,
    },
    groupplot ylabel/.initial={},
    every groupplot y label/.style={
            rotate=90,
        at={($({\pgfplots@group@name\space c1r1.north}-|{\pgfplots@group@name\space c1r1.outer
west})!0.5!({\pgfplots@group@name\space c1r\pgfplots@group@rows.south}-|{\pgfplots@group@name\space c1r\pgfplots@group@rows.outer west})$)},
        anchor=south
    },
    execute at end groupplot/.code={%
      \node [/pgfplots/every groupplot x label]
{\pgfkeysvalueof{/pgfplots/groupplot xlabel}};  
      \node [/pgfplots/every groupplot y label] 
{\pgfkeysvalueof{/pgfplots/groupplot ylabel}};  
    }
}

\def\endpgfplots@environment@groupplot{%
    \endpgfplots@environment@opt%
    \pgfkeys{/pgfplots/execute at end groupplot}%
    \endgroup%
}



%%%%%%%%%%%%%%%%%%%%%%%%%%%%%%%%%%%%%%%%%%%%%%%%%
%   chizen line
%%%%%%%%%%%%%%%%%%%%%%%%%%%%%%%%%%%%%%%%%%%%%%%%%
\pgfplotsset{
  chizen line/.style={
 	cycle list name = chizen lines,
        width  = \columnwidth,
        height = 6cm,
        	axis line style = very thin,
        font = \sffamily,
        % X Axis
        %major x tick style = transparent,
        xlabel style={name=xlabel},
	% Y Axis
        ymajorgrids = true,
        scaled y ticks = false,
        separate axis lines,
	axis lines* = left,
        ymajorgrids = true,
        major tick length = -5,
        %Legend
        legend style={at={(xlabel.south)},yshift=-0.5ex, xshift=-2ex, anchor=north, legend columns = 3, draw=none, column sep = 0.25cm},
        	%legend image code/.code={%
            %        \draw[#1,fill] circle (0.075cm);
            %    },
	}
}

%%%%%%%%%%%%%%%%%%%%%%%%%%%%%%%%%%%%%%%%%%%%%%%%%
%   chizen bar / A Clustered Bar Graph
%%%%%%%%%%%%%%%%%%%%%%%%%%%%%%%%%%%%%%%%%%%%%%%%%
\pgfplotsset{
  chizen bar/.style={
	ybar = 0pt,
  	cycle list name = okabe-ito,
        width  = \columnwidth,
        height = 6cm,
        	axis line style = very thin,
       	bar width=25pt,
        font = \sffamily,
        	% X Axis
        	major x tick style = transparent,
        	enlarge x limits = {abs=1.5cm},
        	% No default xlabel/ylabel: set them in each figure. The legend is
        	% anchored below the x label (node `xlabel'), so set an xlabel.
        	xlabel style={name=xlabel},
	% Y Axis
        	ymajorgrids = true,
        	scaled y ticks = false,
        	separate axis lines,
	axis lines* = left,
        	ymajorgrids = true,
        	major tick length = 0,
        %Legend
        	legend style={at={(xlabel.south)},yshift=-0.5ex, xshift=-2ex, anchor=north, legend columns = 3, draw=none, column sep = 0.25cm},
        	legend image code/.code={\draw[#1] rectangle (0.25cm,0.25cm);}
	}
}

%%%%%%%%%%%%%%%%%%%%%%%%%%%%%%%%%%%%%%%%%%%%%%%%%
%   chizen stacked bar / A Stacked Bar Graph
%%%%%%%%%%%%%%%%%%%%%%%%%%%%%%%%%%%%%%%%%%%%%%%%%
\pgfplotsset{
  chizen stacked bar/.style={
  	ybar stacked,
  	cycle list name = okabe-ito,
        width  = \columnwidth,
        height = 6cm,
        	axis line style = very thin,
       	bar width=25pt,
       	font = \sffamily,
        	% X Axis
        	major x tick style = transparent,
        	enlarge x limits = {abs=1.5cm},
        	% No default xlabel/ylabel: set them in each figure. The legend is
        	% anchored below the x label (node `xlabel'), so set an xlabel.
        	xlabel style={name=xlabel},
	% Y Axis
        	ymajorgrids = true,
        	scaled y ticks = false,
        	separate axis lines,
	axis lines* = left,
        	ymajorgrids = true,
        	major tick length = 0,
        %Legend
        	legend style={at={(xlabel.south)},yshift=-0.5ex, xshift=-2ex, anchor=north, legend columns = 3, draw=none, column sep = 0.25cm},
        	legend image code/.code={\draw[#1] rectangle (0.25cm,0.25cm);}
	}
}


%%%%%%%%%%%%%%%%%%%%%%%%%%%%%%%%%%%%%%%%%%%%%%%%%
%   chizen Likert / Divergent Bar
%%%%%%%%%%%%%%%%%%%%%%%%%%%%%%%%%%%%%%%%%%%%%%%%%
\pgfplotsset{
  chizen Likert/.style={
 	xbar stacked,
        	axis line style = very thin,
       	bar width=5pt,
        	% X Axis
        	%major x tick style = transparent,
        	% No default xlabel: set it in each figure (the legend is anchored below it).
        	xlabel style={name=xlabel},
	xtick = {},
	font = \sffamily,
	% Y Axis
	        %axis y line=none,
	        major y tick style = transparent,
        	%ymajorgrids = true,
        	scaled y ticks = false,
        	separate axis lines,
	axis lines* = left,
        	major tick length = -2,
        %Legend
        	legend style={at={(xlabel.south)},yshift=-2ex, xshift=-2ex, anchor=north, legend columns = 3, draw=none, column sep = 0.25cm},
        	legend image code/.code={\draw[#1] rectangle (0.25cm,0.25cm);}
	}
}

%%%%%%%%%%%%%%%%%%%%%%%%%%%%%%%%%%%%%%%%%%%%%%%%%
%   Tufte Stacked Panel 
%%%%%%%%%%%%%%%%%%%%%%%%%%%%%%%%%%%%%%%%%%%%%%%%%
\pgfplotsset{
  chizen tufte panel/.style={
 ybar, 
  cycle list name = okabe-ito,
	% Generic Bar Graph Options
        width = 1.1\columnwidth,
        height = .325\columnwidth,
        font = \sffamily,
	% X Axis
	enlarge x limits = {abs=.35cm},
	axis line style = very thin,
	% Y Axis
	yticklabel=\empty,
        separate axis lines,
	y axis line style= { draw opacity=0 },
	axis lines* = left,
        axis on top,
        ymajorgrids = true,
        major tick length = 0,
        grid style = {white},
	% Legend Formatting
	title style = {yshift=-.75cm,font=\sffamily\tiny}
	}
}

%%%%%%%%%%%%%%%%%%%%%%%%%%%%%%%%%%%%%%%%%%%%%%%%%
%   Tufte Range Frame
%%%%%%%%%%%%%%%%%%%%%%%%%%%%%%%%%%%%%%%%%%%%%%%%%
\pgfplotsset{
range frame/.style={
    tick align=outside,
    axis line style={opacity=0},
    after end axis/.code={
      \draw ({rel axis cs:0,0}
          -|{axis cs:\pgfplotsdataxmin,0})
        -- ({rel axis cs:0,0}
          -|{axis cs:\pgfplotsdataxmax,0});
      \draw ({rel axis cs:0,0}
          |-{axis cs:0,\pgfplotsdataymin})
        -- ({rel axis cs:0,0}
          |-{axis cs:0,\pgfplotsdataymax});
} }
}

%%%%%%%%%%%%%%%%%%%%%%%%%%%%%%%%%%%%%%%%%%%%%%%%%
%   Tufte Box Plot 
%%%%%%%%%%%%%%%%%%%%%%%%%%%%%%%%%%%%%%%%%%%%%%%%%
\pgfplotsset{
  chizen tufte box/.style={
	% Colour Options
	ybar = 0pt,
	cycle list name = okabe-ito,
        width  = \columnwidth,
        height = 6cm,
        	axis line style = very thin,
       	bar width=1pt,
       	font = \sffamily,
        	% X Axis
        	major x tick style = transparent,
        	enlarge x limits = {abs=1.5cm},
        	xlabel style={name=xlabel},
	% Y Axis
        	ymajorgrids = true,
        	scaled y ticks = false,
        	separate axis lines,
	axis lines* = left,
        	ymajorgrids = true,
        	major tick length = 0,
        %Legend
        %	legend style={at={(xlabel.south)},yshift=2ex, xshift=-2ex, anchor=north, legend columns = 3, draw=none, column sep = 0.25cm},
        %	legend image code/.code={\draw[#1] rectangle (0.25cm,0.25cm);}
 	% Boxplot specific options
	boxplot,
	boxplot/draw direction=y, 
    	%clip = false,
    	every axis plot/.style={
      	mark=o,
      	boxplot/draw direction=y,
      	boxplot/whisker extend=0,
      	boxplot/draw/box/.code={ },
	boxplot/draw/median/.code={%
	 \draw[mark size=2pt,/pgfplots/boxplot/every median/.try]
          \pgfextra
          \pgftransformshift{
            \pgfplotsboxplotpointabbox
              {\pgfplotsboxplotvalue{median}}
              {0.5}
          }
          \pgfsetfillcolor{black}
          \pgfuseplotmark{*}
          \endpgfextra
        	;
      	},
      }
  },
}


%%%%%%%%%%%%%%%%%%%%%%%%%%%%%%%%%%%%%%%%%%%%%%%%%
%   chizen box / Box Plot
%%%%%%%%%%%%%%%%%%%%%%%%%%%%%%%%%%%%%%%%%%%%%%%%%

\pgfplotsset{
    chizen box/.style={
    	% Colour Options
	ybar = 0pt,
	cycle list name = okabe-ito,
        width  = \columnwidth,
        height = 6cm,
        axis line style = very thin,
        font = \sffamily,
        	% X Axis
        	major x tick style = transparent,
        	enlarge x limits = {abs=0.75cm},
        	xlabel style={name=xlabel},
	% Y Axis
        	ymajorgrids = true,
        	scaled y ticks = false,
        	separate axis lines,
	axis lines* = left,
        	ymajorgrids = true,
        	major tick length = 0,
        %Legend
        	legend style={legend columns = 2, draw=none, column sep = 0.25cm},
        	legend image code/.code={\fill[#1] (0,0) circle (0.1cm);}
        % draw whiskers as a single line:
        boxplot/draw/whisker/.code 2 args={%
            \draw[/pgfplots/boxplot/every whisker/.try]
                (boxplot cs:##1) -- (boxplot cs:##2)
            ;
        },%
        %
        % fill the boxes:
        boxplot/every box/.style={
            fill,
        },
        % 
        % the median should be visualized as a thick white line:
        boxplot/every median/.style={
        		thick,
		fill=white, draw=white,
        },
        %boxplot/draw/median/.code={%
        %    \draw[fill=white]
        %        (boxplot cs:\pgfplotsboxplotvalue{median}) circle (3pt)
        %    ;
        %},
        % draw the average as a circle
%        boxplot/average=auto,
%        boxplot/draw/average/.code={%
%	   \draw[fill=white,draw=gray]
%            	(boxplot cs:\pgfplotsboxplotvalue{average}) circle (2pt)
%            ;
%        },
        %
        % do not clip to avoid problems with the median:
        clip=false,
        %
        boxplot/draw direction=y,
        boxplot/whisker extend=0,
        boxplot/every whisker/.style = very thick,
        %
        %
        % width of boxes:
        boxplot/box extend=0.15,
    },
    %
    %
    rshift/.style={
        xshift=\pgfkeysvalueof{/pgfplots/rshift scale},
        legend image post style={xshift=-\pgfkeysvalueof{/pgfplots/rshift scale}},
    },
    lshift/.style={
        xshift=-\pgfkeysvalueof{/pgfplots/lshift scale},
        legend image post style={xshift=\pgfkeysvalueof{/pgfplots/lshift scale}},
    },
    rshift scale/.initial=1em,
    lshift scale/.initial=1em,
}

%%%%%%%%%%%%%%%%%%%%%%%%%%%%%%%%%%%%%%%%%%%%%%%%%
%   OLD NAMES (aliases) -- keep these: existing papers and the manual use them
%%%%%%%%%%%%%%%%%%%%%%%%%%%%%%%%%%%%%%%%%%%%%%%%%
% Each `IDEA ...' style is the `chizen ...' style under its old name. Use
% the chizen names in new figures. Don't remove an alias without checking
% that nothing in the manual or the starter paper still uses it.
\pgfplotsset{
  IDEA line/.style={chizen line},
  IDEA bar/.style={chizen bar},
  IDEA stacked bar/.style={chizen stacked bar},
  IDEA Likert/.style={chizen Likert},
  IDEA tufte panel/.style={chizen tufte panel},
  IDEA tufte box/.style={chizen tufte box},
  IDEA box/.style={chizen box},
}
% A cycle list cannot be aliased by name, so `IDEA lines' repeats the
% entries of `chizen lines'. Keep the two lists identical.
\pgfplotscreateplotcyclelist{IDEA lines}{
{OkabeItoBlue, solid, mark=*, mark options={solid}},
{OkabeItoVermillion, dashed, mark=square*, mark options={solid}},
{OkabeItoBluishGreen, dotted, mark=triangle*, mark options={solid}},
{OkabeItoReddishPurple, dashdotted, mark=diamond*, mark options={solid}},
{OkabeItoOrange, densely dashed, mark=o, mark options={solid}},
{OkabeItoSkyBlue, densely dotted, mark=square, mark options={solid}},
{OkabeItoBlack, loosely dashed, mark=triangle, mark options={solid}},
}

%%%%%%%%%%%%%%%%%%%%%%%%%%%%%%%%%%%%%%%%%%%%%%%%%
% allows multiple lines for the same Y variable in stacked plots
%%%%%%%%%%%%%%%%%%%%%%%%%%%%%%%%%%%%%%%%%%%%%%%%%
\newcommand\resetstackedplots{
\makeatletter
\pgfplots@stacked@isfirstplottrue
\makeatother
}

%%%%%%%%%%%%%%%%%%%%%%%%%%%%%%%%%%%%%%%%%%%%%%%%%
% EOF
%%%%%%%%%%%%%%%%%%%%%%%%%%%%%%%%%%%%%%%%%%%%%%%%%
\makeatother

% Under LuaLaTeX, pgfplots hands some computations to Lua. That Lua backend
% currently fails on some box plots ("linear map received invalid input
% domain inf:-inf"), so use the TeX code path, as pdfLaTeX does.
\pgfplotsset{lua backend=false}

% This caches the PGF Plots stuff to speed up compile times
%\usetikzlibrary{external}
%\tikzexternalize[prefix=submission/tikz/]


%
%   Tables
%
%appendix tables
\usepackage{longtable}
\usepackage{float}
\usepackage{multicol}
\usepackage{colortbl}


% Katja's dashed lines in table stuff
% arydshln rewrites LaTeX's table code and breaks the table structure in a
% tagged PDF (rows get the wrong number of cells, which fails PDF/UA), see
% https://latex3.github.io/tagging-project/tagging-status/ . So it is only
% loaded when tagging is off; with tagging on, \cdashlinelr{2-4} draws a
% solid booktabs \cmidrule instead, and \hdashline/\cdashline are not
% available.
\makeatletter
\IfDocumentMetadataTF{%
\newcommand{\cdashlinelr}[1]{\cmidrule(lr){#1}}
}{%
\usepackage{arydshln}
\def\adl@drawiv#1#2#3{%
	\hskip.5\tabcolsep
	\xleaders#3{#2.5\@tempdimb #1{1}#2.5\@tempdimb}%
	#2\z@ plus1fil minus1fil\relax
	\hskip.5\tabcolsep}
\newcommand{\cdashlinelr}[1]{%
	\noalign{\vskip\aboverulesep
		\global\let\@dashdrawstore\adl@draw
		\global\let\adl@draw\adl@drawiv}
	\cdashline{#1}
	\noalign{\global\let\adl@draw\@dashdrawstore
		\vskip\belowrulesep}}
}
\makeatother


%
%   Stats
%
% acmart already loads etoolbox (ACM's TAPS asks papers not to load it again).
\ifchizenfinal\else
\usepackage{etoolbox}
\fi
\usepackage{siunitx}

\sisetup{
    % range-phrase = --, % make ranges use en-dash instead of "to" (e.g., ages)
    round-mode = places,
    round-precision = 1, % can change precision globally
    round-half = even
}

% Sizes
\newcommand{\by}{\,$\times$\,}
\newcommand{\size}[3][]{\SIrange[range-phrase=\by,%
                                 range-units=single]{#2}{#3}{#1}}
\newcommand{\sizen}[2][]{\SIlist[list-separator=\by,%
                                 list-final-separator=\by,%
                                 list-pair-separator=\by,%
                                %  list-units=single% % repeat (default),single,brackets
                                 ]{#2}{#1}}


% Summary Statistics
\newcommand{\sd}[2][]{$SD=\SI{#2}{#1}$}
\newcommand{\mean}[2][]{$M=\SI{#2}{#1}$}
\newcommand{\mdn}[2][]{$Mdn=\SI{#2}{#1}$}
\newcommand{\ages}[2]{Ages \numrange{#1}{#2}}
\newcommand{\gender}[2]{\num{#1} identified as #2}
\newcommand{\genderlist}[2]{\numlist{#1} identified as #2, respectively}


% Statistical Tests
\newcommand{\fprecision}{1}
\newcommand{\chisqprecision}{1}
\newcommand{\formatprob}[2]{%
    \ifdimless{#2 pt}{.001 pt}{$#1 < .001$}{%
        \ifdimless{#2 pt}{.01 pt}{$#1 < .01$}{%
            $#1 = \num[add-integer-zero=false,%
                      print-zero-integer=false,% siunitx v3: add-integer-zero alone no longer drops the 0
                      round-mode=places,%
                      round-precision=2,%
                      round-half=even]{#2}$%
        }%
    }%
}
% Effect sizes (e.g., partial eta squared) are bounded by 1: two decimals,
% no leading zero, and never "< .001" (that rule is for p-values only).
\newcommand{\formateffect}[2]{%
    $#1 = \num[add-integer-zero=false,%
               print-zero-integer=false,% siunitx v3: add-integer-zero alone no longer drops the 0
               round-mode=places,%
               round-precision=2,%
               round-half=even]{#2}$%
}
\newcommand{\p}[1]{\formatprob{p}{#1}}
\newcommand{\etasqp}[1]{\formateffect{\eta_{p}^2}{#1}}
\newcommand{\F}[3]{{$F_{#1,#2}=\num[round-mode=places,%
                                    round-precision=\fprecision]{#3}$}}
\newcommand{\owF}[2]{{$F_{#1}=\num[round-mode=places,%
                                   round-precision=\fprecision]{#2}$}}
\newcommand{\chisq}[2]{$\chi^2(#1)=\num[round-mode=places,%
                                        round-precision=\chisqprecision]{#2}$}
                                        
\newcommand{\msd}[3][]{\mean[#1]{#2}, \sd[#1]{#3}}
\newcommand{\anova}[4]{\F{#1}{#2}{#3}, \p{#4}}
\newcommand{\anovae}[5]{\F{#1}{#2}{#3}, \p{#4}, \etasqp{#5}}
\newcommand{\owanova}[4]{\owF{#1}{#2}, \p{#3}, \etasqp{#4}}
\newcommand{\chisquare}[3]{\chisq{#1}{#2}, \p{#3}}


% References: biblatex with ACM's styles (acmnumeric, or acmauthoryear), run with biber.
\usepackage[datamodel=acmdatamodel, style=acmnumeric, backend=biber]{biblatex}
\addbibresource{submission/bibliography.bib}


%%
%% \BibTeX command to typeset BibTeX logo in the docs
\AtBeginDocument{%
  \providecommand\BibTeX{{%
    Bib\TeX}}}

%% Rights management information.  This information is sent to you
%% when you complete the rights form.  These commands have SAMPLE
%% values in them; it is your responsibility as an author to replace
%% the commands and values with those provided to you when you
%% complete the rights form.
% \setcopyright{acmlicensed}
% \copyrightyear{2027}
% \acmYear{2027}
% \acmDOI{XXXXXXX.XXXXXXX}

%% These commands are for a PROCEEDINGS abstract or paper. Copy the exact
%% values from your rights confirmation email.
% \acmConference[Conference acronym 'XX]{Make sure to enter the correct
%   conference title from your rights confirmation email}{Month DD--DD,
%   YYYY}{City, Country}
% \acmISBN{978-1-4503-XXXX-X/YYYY/MM}

%%
%% These commands are for a JOURNAL article (e.g., PACM HCI).
% \acmJournal{PACMHCI}
% \acmVolume{0}
% \acmNumber{0}
% \acmArticle{000}
% \acmMonth{0}

%%
%% Submission ID.
%% Use this when submitting an article to a sponsored event. You'll
%% receive a unique submission ID from the organizers
%% of the event, and this ID should be used as the parameter to this command.
%%\acmSubmissionID{123-A56-BU3}

%%
%% The majority of ACM publications use numbered citations and
%% references.  The command \citestyle{acmauthoryear} switches to the
%% "author year" style.
%%\citestyle{acmauthoryear}


%%
%% end of the preamble, start of the body of the document source.
\begin{document}

%%
%% The "title" command has an optional parameter,
%% allowing the author to define a "short title" to be used in page headers.
%% The title is also written into the PDF metadata (dc:title), which screen
%% readers announce, so make it the real title of your paper.
\title[My New CHI Paper]{My New CHI Paper}

%%
%% The "author" command and its associated commands are used to define
%% the authors and their affiliations. With the `anonymous' option the
%% names are replaced automatically, so fill them in now.
\author{Leave Authors Anonymous}
\affiliation{%
  \institution{University of Waterloo}
  \city{Waterloo}
  \state{Ontario}
  \country{Canada}
}
%\email{you@uwaterloo.ca}
%\orcid{0000-0000-0000-0000}

%%
%% By default, the full list of authors will be used in the page
%% headers. Often, this list is too long, and will overlap
%% other information printed in the page headers. This command allows
%% the author to define a more concise list
%% of authors' names for this purpose.
%\renewcommand{\shortauthors}{Author1 et al.}



%%
%% The abstract is a short summary of the work to be presented in the
%% article.



\begin{abstract}
\begin{figure}[H]
    \centering
    \includegraphics[width=.32\textwidth]{submission/figures/chi2021.png} 
    \includegraphics[width=.32\textwidth]{submission/figures/university-of-waterloo-vertical-logo.png} 
    \includegraphics[width=.32\textwidth]{submission/figures/uwhci-full-cs.png}
    \caption{Example teaser figure.}
    \label{fig:teaser}
\end{figure}

\end{abstract}

%%
%% The code below is generated by the tool at http://dl.acm.org/ccs.cfm.
%% Please copy and paste the code instead of the example below.
%%
% \begin{CCSXML}
% <ccs2012>
% <concept>
% <concept_id>10003120.10003130.10011762</concept_id>
% <concept_desc>Human-centered computing~Empirical studies in collaborative and social computing</concept_desc>
% <concept_significance>500</concept_significance>
% </concept>
% </ccs2012>
% \end{CCSXML}

% \ccsdesc[500]{Human-centered computing~Empirical studies in collaborative and social computing}

%%
%% Keywords. The author(s) should pick words that accurately describe
%% the work being presented. Separate the keywords with commas.
%\keywords{keyword one, keyword two, keyword three}


%%
%% This command processes the author and affiliation and title
%% information and builds the first part of the formatted document.
\maketitle


\begin{markdown}
Hello, World!
\end{markdown}

%
% Add files as below for each section, or just start writing here
%
% Every figure needs alt text: put \Description{...} in every figure
% environment (see the CHI Zen manual, Section "Accessibility").
%

%
\vspace{2em}
\begin{quote}
    \emph{Great things are done by a series of small things brought together. --- Vincent Van Gough}
\end{quote}


%%%%%%%%%%%%%%%%%%%%%%%%%%%%%%%%%%%%%%%%%%%%%%%%%%%%%%%%%%%%%%%%%%%%%%%%%%%%%
\section{Introduction}
%%%%%%%%%%%%%%%%%%%%%%%%%%%%%%%%%%%%%%%%%%%%%%%%%%%%%%%%%%%%%%%%%%%%%%%%%%%%%

Even the most experienced researchers are constantly looking for ways to improve their papers. And there are many different aspects of a paper to improve on: from high-level concerns like experimental design, positioning the paper within the literature, and consideration of ethical issues, to sweating the details in a statistical analyis, or getting that graph \emph{just right} to show off your earth-shattering results. But sometimes it's hard to put all the pieces together. That's where this project comes in. 

This project --- CHI Zen, a play on the words `KaiZen' (\begin{CJK}{UTF8}{min}改善\end{CJK}) or `continuous improvement' --- describes common pitfalls and best practices based on my experience writing papers for a number of years. That experience reflects some of my own preferences, but also a lot that I've learned from working with other HCI researchers. It also contains a \epmh{lot} of examples that you are free to copy, modify, or re-use as often as you'd like. The goal is to have something of an \emph{information zoo} \cite{heer2010tour} that can help you think about how to best visualize results, and hopefully make the process of creating that visualization easy, too. 

This document, also provides a sample reference for ``IDEA Plots,'' a set of templates that I developed for \texttt{PGFPlots}. Please use IDEA Plots if you find it useful. If you do, I'd love to hear about it, and if you have any ways to imrpove the templates. You might also prefer to use a different graphing package, and that's cool too. The examples contained in this document hopefully can provide tips or shortcuts for use in your own papers. I suspect that you can generate more or less the same graphs with your tool of choice. 

Finally, my intention in making this project available via GitHub is to allow it to grow. If you find it useful, please drop me a line. If you find a way to improve it, please push the changes to GitHub. It might also just be an important prompt to have a conversation about best practices, or an opportunity to explore different ways of thinking about or presenting research. In any case, I hope that it's an opportunity to think about how we can improve our research practices, one small step at a time. 



%
\vspace{2em}
\begin{center}
    \emph{Excellence is not a destination, it is a continuous journey that never ends.  --- Brian Tracy}
\end{center}

%%%%%%%%%%%%%%%%%%%%%%%%%%%%%%%%%%%%%%%%%%%%%%%%%%%%%%%%%%%%%%%%%%%%%%%%%%%%%
\section{\LaTeX\ and IDEA Plots}
%%%%%%%%%%%%%%%%%%%%%%%%%%%%%%%%%%%%%%%%%%%%%%%%%%%%%%%%%%%%%%%%%%%%%%%%%%%%%

\LaTeX is far from a perfect writing tool. For one thing, it's a `typesetting environment' and not a word processor, so little things like the difference between a ` and a ' can end up being important. When you're writing a paper, you generally want to focus on the words you're writing, and not getting the syntax right. So why use \LaTeX?

\begin{enumerate}
    \item Most formatting is done `for free' by the SIG CHI template, and it does a great job of this for you
    \item The University of Waterloo provides a Pro Overleaf account (think Google Docs for academic writing) if you login with your UW userid
    \item \LaTeX\ integrates well with other (free) tools like R, PGFPLots, etc. 
    \item Papers get rejected all the time, and reformatting is very fast and easy
\end{enumerate}

Since \LaTeX\ can be complex, this project also embodies some changes to the template that make our typical workflow a little easier:

\begin{enumerate}
    \item Sections are separated out into files, so that authors can edit each section without interrupting each other
    \item Additional packages:
    \item IDEA Plots --- see below. 
    \item ... 
\end{enumerate}



\subsection{PGFPlots and IDEA Plots}
As as I was completeing my PhD, I decided to spend some time trying to learn to make better figures since this was something I expected to spend a lot of time doing over the rest of my career, and I might as well think about how to do it well from the start. Some of the things I'd hoped to improve on were using scalable graphics, having better control over the width/height of figures, and being able to use experimental data itself to create the images (better fit with a typical academic workflow). I also prefer a minimalist feel, and wanted my papers to reflect that (And even if you prefer something different, this is a good starting point, you can always make things more complex where needed). 

So after doing some research, I found that \texttt{PGFPlots} was a good fit for my needs. But it wasn't perfect, and in particular can be complicated to get right. So in the interest of saving the time of re-inventing the wheel, and (hopefully) enabling others to use my work for their own papers, I decided to create a set of templates that would encapsulate the common types of graphs I use in HCI research. I called these templates \texttt{IDEA Plots}. As I developed the templates, they grew to include style choices, colour palettes, and some common shortcuts for layout. I expect they'll continue to evolve for a while to come.  

All of the examples in this document use the templates from \texttt{IDEA Plots}, with some in-place customization based on the needs of the graph. If you'd like to use those templates in your own work, be sure to include \texttt{IDEAPlots.tex} in your own project, and to \texttt{\textbackslash include\{IDEAPlots.tex\}} early on in your paper.  

%

%%%%%%%%%%%%%%%%%%%%%%%%%%%%%%%%%%%%%%%%%%%%%%%%%%%%%%%%%%%%%%%%%%%%%%%%%%%%%
%\section{GAME DESIGN}
%%%%%%%%%%%%%%%%%%%%%%%%%%%%%%%%%%%%%%%%%%%%%%%%%%%%%%%%%%%%%%%%%%%%%%%%%%%%%



%

\newpage
\vspace{2em}
\begin{center}
    \emph{If you define the problem correctly, you almost have the solution. --- Steve Jobs}
\end{center}

%%%%%%%%%%%%%%%%%%%%%%%%%%%%%%%%%%%%%%%%%%%%%%%%%%%%%%
\section{Tables}
%%%%%%%%%%%%%%%%%%%%%%%%%%%%%%%%%%%%%%%%%%%%%%%%%%%%%%
As a rule of thumb, if you can show your data with a table, consider doing that first. Tables are quite effective for showing real data, and usually easy to interpret. The Vis literature even says that tables might be more convincing for a skeptical audience \cite{pandey2014persuasive}!

In general, I suggest the following guidelines for formatting: 
\begin{enumerate}
    \item Remove vertical lines
    \item Only use horizontal lines where necessary
    \item Use \texttt{ \textbackslash toprule}, \texttt{ \textbackslash midrule}, and \texttt{ \textbackslash bottomrule} to frame your table
\end{enumerate}

You can see an example in \autoref{tab:exampletable}, below.

\begin{center}
\begin{table*}[htbp!]
\small
\begin{center}
\begin{tabular}{l l l l l}
\toprule
Condition & \multicolumn{4}{l}{Teamwork Measures}\\
\cmidrule(r){2-5}
& \multicolumn{2}{l}{Measure 1} & Measure 2 & Measure 3 \\
\cmidrule(r){2-3} 
&  Sub 1 & Sub 2 &  \\ 


\midrule


Condition 1 & .478 & .468 & 1.68 & 12.3  \\
                  & (.0259) & (.031) & (.703) & (5.93) \\
\\

Condition 2 & .444 & .451 & 1.32 & 18.4 \\
                              & (.0478) & (.0585) & (.319) & (2.76) \\
\\

Condition 3 &  & .410 &  2.857 &  \\
                     &        & (.0384)  & (.350) & \\
\\


ANOVA Results & $F_{(1,12)} = 2.27,$ & $F_{(2,18)} = 2.76,$  & $F_{(2,18)} = 16.16,$ & $F_{(1,12)} = 6.16,$ \\
                          & $p = .158$ & $p = .090  $ &  $p < .0001$* &  $p = 0.029$ * \\
\bottomrule
\end{tabular}

\caption[a]{Mean values and standard deviations (in parentheses) for teamwork measures, and ANOVA results for comparisons between experimental conditions. Significant results denoted by *.}
\label{tab:exampletable}
\end{center}
\end{table*}
\end{center}

%\newpage



%%%%%%%%%%%%%%%%%%%%%%%%%%%%%%%%%%%%%%%%%%%%%%%%%%%%%%%%%%%%%%%%%%%%%%%%%%%%%
\section{Showing Your Data}
%%%%%%%%%%%%%%%%%%%%%%%%%%%%%%%%%%%%%%%%%%%%%%%%%%%%%%%%%%%%%%%%%%%%%%%%%%%%%


- Figures and Tables in your results section often provide reviewers a first look at your results
- can also make the often dense results section more accessible to readers who are unfamiliar with your research area
- for readers with specific experience, particularly in replicated experimental settings such as a Fitts's Law study, can be the most important summary of research results in the paper 


Some rules of thumb from Munzer \cite{munznervisualization}:
\begin{enumerate}
	\item No unjustified 3D
	\item No unjustified 2D
	\item Eyes beat memory
	\item Get it right in Black and White
	\item Function first, form next
\end{enumerate}


Munzner \cite{munznervisualization}, and in particular Chapter 5, is a useful resource for determining which Marks and Channels to use in a graph, and lists them in order of effectiveness

- the style suggested in this document is informed and inspired by Tufte \cite{tufte1983visual}, who emphasizes minimalist presentation of quantitative data. Tufte's principles of information visualization 

 
 
%  For example, some good examples:
% Bartneck and Hu \cite{Bartneck:2009:SAC:1518701.1518810}


% Senellart \cite{Senellart:2013:DYR:2513166.2514938} discusses some useful cases of `bad' style, which superficially may appear to be picky but offer important guidance towards submitting polished papers. 
 
 
% Munzer \cite{Munzner:2008:PPW:1422919.1422927} 
% - some good advice in general, but more focused on info vis research

% - text is large and legible
% - all axes are labelled, including units of measure
% - axes start at 0, unless explicitly justified (for a good reason!)













%%%%%%%%%%%%%%%%%%%%%%%%%%%%%%%%%%%%%%%%%%%%%%%%%%%%%%%%%%
\subsection{Comparing Performance: Bar Charts and Box Plots}
%%%%%%%%%%%%%%%%%%%%%%%%%%%%%%%%%%%%%%%%%%%%%%%%%%%%%%%%%%

- common to compare performance, efficiency, error rate, or other dependent measures between two study groups. For example, our new technique for 3D pointing, compared to a computer mouse. 

- most simple way of doing so is to show the means and standard error in a bar graph, and this is commonly accepted throughout the ACM. 
- For example, Figure \ref{bargraph} is an example drawn from a UIST publication \cite{REF}.


\begin{figure}
\begin{tikzpicture}
\begin{axis}[IDEA bar,
        symbolic x coords={non-occluded,occluded},
        xtick = {non-occluded,occluded},
        enlarge x limits=0.5,
	    ymin = 0, ymax = 10,
        ylabel = {Selection Time (s)},
        width=\columnwidth,
        height = 5cm,
]

%Depth Ray
\addplot[style={draw=none,fill=AHSLight},error bars/.cd, y dir = both, y explicit, error bar style={black}]
coordinates {(occluded, 6.8)  +- (1.7,1.7)};
\addlegendentry{Depth Cursor}

%Tiltcasting
\addplot+[style={draw=none, fill=SchoolRedLight,line width = 0pt},error bars/.cd, y dir = both, y explicit, error bar style={black}]
coordinates {(non-occluded,3.24) +- (.712,.712)
		    (occluded,4)  +- (.922,.922)};
\addlegendentry{Tiltcasting}

%Smartcasting
\addplot+[style={draw=none, fill=EnvironmentLight},error bars/.cd, y dir = both, y explicit, error bar style={black}]
coordinates {(non-occluded, 2.10)  +- (.42,.42)};
\addlegendentry{Smartcasting}

\end{axis}
\end{tikzpicture}
\caption{An example bar graph with error bars. Data is specified in the latex file.}
\label{bargraph}
\end{figure}


- to be more accurate and descriptive, box plots are often used
- similar to a bar graph with error bars, but also indicate quartiles and outliers 
- typically, better practice to include a box plot when possible -- but can be substantially more difficult to create in Microsoft Word. 




%%%%%%%%%%%%%%%%%%%%%%%%%%%%%%%%%%%%%%%%%%%%%%%%%%%%%%%%%%
\subsection{Questionnaire Responses}
%%%%%%%%%%%%%%%%%%%%%%%%%%%%%%%%%%%%%%%%%%%%%%%%%%%%%%%%%%

- Common to want to show questionnaire respones, particularly to, e.g., Likert scale questions
- Jim's preference: stacked divergent bar graph

\begin{figure}
\begin{tikzpicture}
\pgfplotstableread[col sep = comma]{submission/data/Cleric_AvatarIdentificationStacked.csv}\ClericIdentificationData
\pgfplotstableread[col sep = comma]{submission/data/Monster_AvatarIdentificationStacked.csv}\MonsterIdentificationData

\begin{axis}[
    IDEA Likert,
    height = 6.5cm,
    width = .75\columnwidth,
    % Y AXIS
    symbolic y coords = {Ava_Conn,	Ava_WE,	AvavsPhys,	AvaVsPers,	AvaVsPhysIdeal,	AvaVsPersIdeal},
    ytick={Ava_Conn, Ava_WE, AvavsPhys, AvaVsPers, AvaVsPhysIdeal, AvaVsPersIdeal},
    yticklabels={Connectedness, Refer to \\ Avatar as ``We'', Physical \\ Similarity, Personality\\ Similarity, Physical Ideal, Ideal Personality},
    yticklabel style={align=right,font={\sffamily\small}},
    y axis line style={draw=none},
    y dir = reverse,
    ymajorgrids = false,
    % X AXIS
    xmin=-100, xmax=100,
    xlabel={\% of Responses},
    xlabel style = {font=\small\sffamily},
    xticklabel style = {font=\small\sffamily},
    xtick = {-100,-50,0,50,100},
    xticklabels = {100, 50, 0, 50, 100},
    extra x ticks = {-75, 75},
    extra x tick labels = {$\leftarrow$ Disagree, Agree $\rightarrow$},
    every extra x tick/.style={major tick length=0,yshift={-8pt}},
    ]
    \addplot[draw=none,fill=none, forget plot] coordinates {(0,Ava_Conn)(0,Ava_WE)(0,AvavsPhys)(0,AvaVsPers)(0,AvaVsPhysIdeal)(0,AvaVsPersIdeal)};
    \addplot[draw=none,fill=Cleric_3, forget plot, bar shift=.1cm] table [x expr={\thisrow{4}*-0.5},y=Measure] {\ClericIdentificationData};
    \addplot[draw=none,fill=Cleric_2, forget plot, bar shift=.1cm] table [x expr={\thisrow{3}*-1},y=Measure] {\ClericIdentificationData};
    \addplot[draw=none,fill=Cleric_1, forget plot, bar shift=.1cm] table [x expr={\thisrow{2}*-1}, y=Measure ] {\ClericIdentificationData};
    \addplot[draw=none,fill=Cleric_0, forget plot, bar shift=.1cm] table [x expr={\thisrow{1}*-1}, y=Measure ] {\ClericIdentificationData};
    \resetstackedplots
    \addplot[draw=none,fill=none, forget plot] coordinates {(0,Ava_Conn)(0,Ava_WE)(0,AvavsPhys)(0,AvaVsPers)(0,AvaVsPhysIdeal)(0,AvaVsPersIdeal)};
    \addplot[draw=none,fill=Monster_3, forget plot, bar shift=-.1cm] table [x expr={\thisrow{4}*-0.5},y=Measure] {\MonsterIdentificationData};
    \addplot[draw=none,fill=Monster_2, forget plot, bar shift=-.1cm] table [x expr={\thisrow{3}*-1},y=Measure] {\MonsterIdentificationData};
    \addplot[draw=none,fill=Monster_1, forget plot, bar shift=-.1cm] table [x expr={\thisrow{2}*-1}, y=Measure ] {\MonsterIdentificationData};
    \addplot[draw=none,fill=Monster_0, forget plot, bar shift=-.1cm] table [x expr={\thisrow{1}*-1}, y=Measure ] {\MonsterIdentificationData};
\end{axis}

\begin{axis}[
    IDEA Likert,
    height = 6.5cm,
    width = .75\columnwidth,
    % Y AXIS
    symbolic y coords = {Ava_Conn,	Ava_WE,	AvavsPhys,	AvaVsPers,	AvaVsPhysIdeal,	AvaVsPersIdeal},
    ytick={Ava_Conn, Ava_WE, AvavsPhys, AvaVsPers, AvaVsPhysIdeal, AvaVsPersIdeal},
    yticklabels={Connectedness, Refer to \\ Avatar as ``We'', Physical \\ Similarity, Personality\\ Similarity, Physical Ideal, Ideal Personality},
    yticklabel style={align=right,font={\sffamily\small}},
    y axis line style={draw=none},
    y dir = reverse,
    ymajorgrids = false,
    % X AXIS
    xmin=-100, xmax=100,
    xlabel={\% of Responses},
    xlabel style = {font=\small\sffamily},
    xticklabel style = {font=\small\sffamily},
    xtick = {-100,-50,0,50,100},
    xticklabels = {100, 50, 0, 50, 100},
    extra x ticks = {-75, 75},
    extra x tick labels = {$\leftarrow$ Disagree, Agree $\rightarrow$ },
    every extra x tick/.style={major tick length=0,yshift={-8pt}},
    ]
    \addplot[draw=none,fill=none, forget plot] coordinates {(0,Ava_Conn)(0,Ava_WE)(0,AvavsPhys)(0,AvaVsPers)(0,AvaVsPhysIdeal)(0,AvaVsPersIdeal)};
    
    \addplot[draw=none,fill=Cleric_3, forget plot, bar shift=.1cm] table [x expr={\thisrow{4}*0.5}, y=Measure ] {\ClericIdentificationData};
    \addplot[draw=none,fill=Cleric_4, forget plot, bar shift=.1cm] table [x expr={\thisrow{5}}, y=Measure ] {\ClericIdentificationData};
    \addplot[draw=none,fill=Cleric_5, forget plot, bar shift=.1cm] table [x expr={\thisrow{6}}, y=Measure ] {\ClericIdentificationData};
    \addplot[draw=none,fill=Cleric_6, forget plot, bar shift=.1cm] table [x expr={\thisrow{7}}, y=Measure ] {\ClericIdentificationData}; \label{plot:cleric6}
    \resetstackedplots
    \addplot[draw=none,fill=none, forget plot] coordinates {(0,Ava_Conn)(0,Ava_WE)(0,AvavsPhys)(0,AvaVsPers)(0,AvaVsPhysIdeal)(0,AvaVsPersIdeal)};
    \addplot[draw=none,fill=Monster_3, forget plot, bar shift=-.1cm] table [x expr={\thisrow{4}*0.5}, y=Measure ] {\MonsterIdentificationData};
    \addplot[draw=none,fill=Monster_4, forget plot, bar shift=-.1cm] table [x expr={\thisrow{5}}, y=Measure ] {\MonsterIdentificationData};
    \addplot[draw=none,fill=Monster_5, forget plot, bar shift=-.1cm] table [x expr={\thisrow{6}}, y=Measure ] {\MonsterIdentificationData};
    \addplot[draw=none,fill=Monster_6, forget plot, bar shift=-.1cm] table [x expr={\thisrow{7}}, y=Measure ] {\MonsterIdentificationData};
    %\addlegendimage{only marks, mark=o}
    %\addlegendimage{only marks, mark=o}
    %\legend{Neutral, Agree, Strongly Agree, Disagree,Strongly Disagree}
    after end axis/.code={
        \node at (axis cs:80,Ava_Conn) [anchor=east, ,font=\tiny] {\AsteriskBold};
        \node at (axis cs:80,Ava_WE) [anchor=east, font=\tiny] {\AsteriskBold};
        %\draw [white] (axis cs:0,AvaVsPersIdeal) -- (axis cs:0,Ava_Conn);
        %\draw [white] (axis cs:-50,AvaVsPersIdeal) -- (axis cs:-50,Ava_Conn);
        %\draw [white] (axis cs:50,AvaVsPersIdeal) -- (axis cs:50,Ava_Conn);
    }
    \coordinate (legend) at (axis description cs:1.15,.85);
\end{axis}

% this is a dummy `axis' environment only to create the legend
\matrix [
    matrix of nodes, 
    every node/.style={anchor=center}, 
    ] at (legend) {
        |[fill=Cleric_0]| & |[fill=Cleric_1]| & |[fill=Cleric_2]| & |[fill=Cleric_3]| & |[fill=Cleric_4]| & |[fill=Cleric_5]| & |[fill=Cleric_6]| & |[font=\small\sffamily]|Cleric \\
        |[fill=Monster_0]| & |[fill=Monster_1]| & |[fill=Monster_2]| & |[fill=Monster_3]| & |[fill=Monster_4]| & |[fill=Monster_5]| & |[fill=Monster_6]| & |[font=\small\sffamily]|Monster \\
    };
\end{tikzpicture}
    
\caption{Summary of avatar identification responses, based on 7-point Likert scales. Mann Whitney U tests revealed significant differences between \textsc{Avatar} groups for `Connectness' and `Refer to Avatar as We'. }
\label{fig:avatar_identification}
\end{figure}




%%%%%%%%%%%%%%%%%%%%%%%%%%%%%%%%%%%%%%%%%%%%%%%%%%%%%%%%%%
\subsubsection{Showing Trends over Time}
%%%%%%%%%%%%%%%%%%%%%%%%%%%%%%%%%%%%%%%%%%%%%%%%%%%%%%%%%%



\begin{figure}[t!]
\begin{tikzpicture}[baseline]
\begin{axis}[IDEA line,
	% Graph Options
	width=.475*\columnwidth,
	height=4cm,
	title=\texttt{r/stop\-drinking},
	% X Axis Options
	%date coordinates in=x,
	%date ZERO=2014-01-01,
    %xticklabel=\year,
    %xmin=2014-01-01,
    %xmax=2018-01-01,
    xlabel=Year,
	xtick={1,13,25,37,49},                              %% Jim: apologies... this is a hack
    xticklabels={2014,2015,2016,2017,2018},             %%      will figure out what's going on here if we have time later
	%xtick ={1,2,3,4,5,6,7,8,9},
	% Y Axis Options
	ylabel = Distinct Users / Active Threads,
	ylabel style = {font=\tiny},
	ymin = 0, ymax = 6000,
	ytick = {0,2000,4000,6000},
	ytick style = {font=\tiny},
    %legend image code/.code={\draw[#1] circle (0.1cm);}
]

\pgfplotstableread[col sep = comma]{submission/data/stopdrinking-usagebymonth.csv}\mydata;
\addplot+[draw=ScienceDark,fill=none,thick] table[x=index,y=users] {\mydata};
\addplot+[draw=ArtsLight,fill=none,thick] table[x=index,y=threads] {\mydata};

\end{axis}
\end{tikzpicture}%
%
\begin{tikzpicture}[baseline]
\begin{axis}[IDEA line,
	% Graph Options
	width=.475*\columnwidth,	
	height=4cm,
	title=\texttt{r/Opiates\-Recovery},
	% X Axis Options
	xtick={1,13,25,37,49},                              %% Jim: apologies... this is a hack
    xticklabels={2014,2015,2016,2017,2018},             %%      will figure out what's going on here if we have time later
	%date coordinates in=x,
	%date ZERO=2014-01-01,
    %xticklabel=\year,
    %xmin=2014-01-01,
    %xmax=2018-01-01,
	xlabel = Year,
	%xtick ={1,2,3,4,5,6,7,8,9},
	% Y Axis Options
	ylabel = Distinct Users / Active Threads,
	ylabel style = {font=\tiny},
	ymin = 0, ymax = 800,
	ytick = {0,200,400,600,800},
	ytick style = {font=\tiny},
	legend style={at={(0.5,-1)},anchor=south}
    %legend image code/.code={\draw[#1] circle (0.1cm);}
]

\pgfplotstableread[col sep = comma]{submission/data/opiaterecovery-usagebymonth.csv}\mydata;
\addplot+[draw=ScienceDark,fill=none,thick] table[x=index,y=users] {\mydata};
\addlegendentry{Distinct Users}
\addplot+[draw=ArtsLight,fill=none,thick] table[x=index,y=threads] {\mydata};
\addlegendentry{Active Threads}

\end{axis}
\end{tikzpicture}
%\setlength{\belowcaptionskip}{-20pt}
\caption{Distinct users and active threads for \texttt{r/stop\-drinking} (top) and \texttt{r/Opiates\-Recovery} (bottom) showing that both counts are experiencing an upwards trend over time for both subreddits.} 

\label{fig:activity}
\end{figure}



\begin{figure*}[tb!]
\begin{center}
\begin{tikzpicture}
\begin{groupplot}[
   group style={
       group size=1 by 5,
       x descriptions at=edge bottom,
       y descriptions at=edge left,
       vertical sep=1pt,
       horizontal sep=0pt},
  IDEA tufte panel, 
	% X Axis
	width=\columnwidth,
	xmin = 1990, xmax = 2015,
	xtick = {1990,1994,1998,2002,2006,2010,2015},
	x tick label style={font=\small,rotate=90, /pgf/number format/1000 sep=},
	axis line style={thick},
	% Y Axis
	ytick={10,20,30,40,50,60,70,80,90},
        ymin = 0,
        ymax = 50,
      ]


% Bibliographic
\nextgroupplot[bar width=7pt, height = 2.25cm, ymax = 25, title={Bibliographic}]
\addplot+[draw=none, fill=black!50] 
 coordinates {(1990,20) (1992,4) (1994,0) (1996,0) (1998,2) (2000,0) (2002,0) (2004,2) (2006,2) (2008,1) (2010,0) (2011,0) (2012,0) (2013,1) (2014,0) (2015,1) };       
  
% Not Empirical      
\nextgroupplot[bar width=7pt, height = 3.5cm, ymax=70, title={Not Empirical}]
\addplot+[draw=none, fill=black!50] 
coordinates {(1990,37) (1992,55) (1994,33) (1996,38) (1998,54) (2000,35) (2002,20) (2004,18) (2006,24) (2008,13) (2010,5) (2011,31) (2012,9) (2013,4) (2014,6) (2015,7) };

% Explanatory
\nextgroupplot[bar width=7pt, height = 3cm, title={Explanatory}]
\addplot+[draw=none, fill=black!50]  
coordinates {(1990,3) (1992,4) (1994,5) (1996,0) (1998,0) (2000,5) (2002,15) (2004,16) (2006,29) (2008,25) (2010,21) (2011,16) (2012,8) (2013,14) (2014,15) (2015,12) };

% D & E
\nextgroupplot[bar width=7pt, ymax=60, height = 3.5cm, title={Design and Evaluation}]
\addplot+[draw=none, fill=blue!50] 
coordinates {(1990,17) (1992,13) (1994,26) (1996,22) (1998,15) (2000,27) (2002,39) (2004,35) (2006,23) (2008,19) (2010,21) (2011,10) (2012,28) (2013,26) (2014,24) (2015,22) };

% Descriptive
\nextgroupplot[bar width=7pt, ymax = 70, height = 3.5cm, title={Descriptive}]
\addplot+[draw=none,fill=black!50]
coordinates {(1990,23) (1992,23) (1994,36) (1996, 40) (1998,29) (2000,32) (2002,27) (2004,29) (2006,23) (2008,42) (2010,53) (2011,42) (2012,55) (2013,55) (2014,55) (2015,58) };


\end{groupplot}
\end{tikzpicture}
\caption{An example of Tufte's stacked panel graph, which presents the same information as a stacked bar graph but better enables comparisons across each vertical dimension. Best practice is to use black and white when possible, with sparing colour for emphasis where appropriate. }
\label{stackedpanel}
\end{center}
\end{figure*}




\begin{figure}
\begin{tikzpicture}
\begin{groupplot}[
   group style={
       group size=1 by 5,
       x descriptions at=edge bottom,
       y descriptions at=edge left,
       vertical sep=10pt,
       horizontal sep=0pt},
  IDEA line,
	% Graph Options
	height=3cm,
	 separate axis lines,
	% X Axis Options
	xmin=0, xmax= 35,
	xtick={1,5,10,15,20,25,30, 35},
	%date coordinates in=x,
	%date ZERO=2014-01-01,
    %xticklabel=\year,
    %xmin=2014-01-01,
    %xmax=2018-01-01,
	xlabel = Number of NPC Requests,
	xlabel style = {font=\small\sffamily},
    xticklabel style = {font=\small\sffamily},
	%xtick ={1,2,3,4,5,6,7,8,9},
	% Y Axis Options
	%ylabel = Proportion of Responses,
	ymajorgrids = false,
	ylabel shift = -.3cm,
	ymin = 0, ymax = 1,
	ytick = {0, 1.00},
	ylabel style = {font=\small\sffamily},
    yticklabel style = {font=\small\sffamily},
	%ytick style = {font=\tiny},
	title style = {yshift=-.75cm,font=\sffamily\small}
]
\pgfplotstableread[col sep = comma]{submission/data/CBT_PropOverTime.csv}\cbtovertimedata;

\nextgroupplot[xtick style = {white}, title style = {yshift=+.35cm,font=\sffamily\small}, title={Not Damaging}]
\addplot+[mark=*, mark size=1.5pt, draw=gray, only marks,fill=gray,thin] table[x=Query,y=NotDamaging] {\cbtovertimedata};
\addplot [solid, draw=black, fill=none,ultra thick] table[y={create col/linear regression={y=NotDamaging}}] {\cbtovertimedata};
%\addlegendentry{Not Damaging}

\nextgroupplot[xtick style = {white}, title={Empathize}]
\addplot+[mark=*, mark size=1.5pt, draw=gray, only marks,fill=gray,thin] table[x=Query,y=Empathize] {\cbtovertimedata};
\addplot [solid, draw=black, fill=none,ultra thick] table[y={create col/linear regression={y=Empathize}}] {\cbtovertimedata};
%\addlegendentry{Empathize}

\nextgroupplot[xtick style = {white}, title style = {yshift=+.25cm,font=\sffamily\small}, title={Reframe}]
\addplot+[mark=*, mark size=1.5pt, draw=gray, only marks,fill=gray,thin] table[x=Query,y=Reframe] {\cbtovertimedata};
\addplot [solid, draw=black, fill=none,ultra thick] table[y={create col/linear regression={y=Reframe}}] {\cbtovertimedata};
%\addlegendentry{Reframe}

\nextgroupplot[xtick style = {white}, title={Encourage}]
\addplot+[mark=*, mark size=1.5pt, draw=gray, only marks,fill=gray,thin] table[x=Query,y=Encourage] {\cbtovertimedata};
\addplot [solid, draw=black, fill=none,ultra thick] table[y={create col/linear regression={y=Encourage}}] {\cbtovertimedata};
%\addlegendentry{Encourage}

\nextgroupplot[title={Solution}]
\addplot+[mark=*, mark size=1.5pt, draw=gray, only marks,fill=gray,thin] table[x=Query,y=Solution] {\cbtovertimedata};
\addplot [solid, draw=black, fill=none,ultra thick] table[y={create col/linear regression={y=Solution}}] {\cbtovertimedata};
%\addlegendentry{Solution}


%\addplot+[mark=x, only marks,draw=Monster_0,fill=none,semithick] table[x=Query,y=Proportion, y error plus=CI-upper, y error minus=CI-lower] {\barbariansolutiondata};
%\addlegendentry{Monster}
%\addplot [draw=Monster_0, fill=none,ultra thick] table[y={create col/linear regression={y=Proportion}}] {\barbariansolutiondata};

\end{groupplot}
\end{tikzpicture}

%\setlength{\belowcaptionskip}{-20pt}
\caption{The proportion of responses that were classified as meeting five criteria: Empathize, Reframe, Encourage, Solution, and Not Damaging. The three categories related to CBT, Empathize, Reframe, and Encourage, were found to decrease over time ($p=0.05$). However, we did not find these changes for Solution and Not Damaging. 
%decreased over time for all participants, however responses by participants in the \textsc{Cleric} group decreased more quickly than those in the \textsc{Monster} group.
} 

\label{fig:SolutionsOverQueryNumber}
\end{figure}
























%%%%%%%%%%%%%%%%%%%%%%%%%%%%%%%%%%%%%%%%%%%%%%%%%%%%%%%%%%%%%%%%%%%%%%%%%%%%%%
%\section{Discussion \& Implications}
%%%%%%%%%%%%%%%%%%%%%%%%%%%%%%%%%%%%%%%%%%%%%%%%%%%%%%%%%%%%%%%%%%%%%%%%%%%%%


%%%%%%%%%%%%%%%%%%%%%%%%%%%%%%%%%%%%%%%%%%%%%%%%%%%%%%%%%%%%%%%%%%%%%%%%%%%%%%
%\section{Limitations}
%%%%%%%%%%%%%%%%%%%%%%%%%%%%%%%%%%%%%%%%%%%%%%%%%%%%%%%%%%%%%%%%%%%%%%%%%%%%%



%%%%%%%%%%%%%%%%%%%%%%%%%%%%%%%%%%%%%%%%%%%%%%%%%%%%%%%%%%%%%%%%%%%%%%%%%%%%%%
%\section{Conclusions}
%%%%%%%%%%%%%%%%%%%%%%%%%%%%%%%%%%%%%%%%%%%%%%%%%%%%%%%%%%%%%%%%%%%%%%%%%%%%%



%%%%%%%%%%%%%%%%%%%%%%%%%%%%%%%%%%%%%%%%%%%%%%%%%%%%%%%%%%%%%%%%%%%%%%%%%%%%%%
\section{Acknowledgements}
%%%%%%%%%%%%%%%%%%%%%%%%%%%%%%%%%%%%%%%%%%%%%%%%%%%%%%%%%%%%%%%%%%%%%%%%%%%%%


\newpage
\printbibliography

%\input{submission/11-appendix.tex}

\end{document}
\endinput
%%
%% End of file.
